% see the original template for more detail about bibliography, tables, etc: https://www.overleaf.com/latex/templates/handout-design-inspired-by-edward-tufte/dtsbhhkvghzz

\documentclass{tufte-handout}

%\geometry{showframe}% for debugging purposes -- displays the margins

\usepackage{amsmath}

\usepackage{hyperref}

\usepackage{fancyhdr}

\usepackage{hanging}

\hypersetup{colorlinks=true,allcolors=[RGB]{97,15,11}}

\fancyfoot[L]{\emph{History of Media Studies}, vol. 6, 2026}


% Set up the images/graphics package
\usepackage{graphicx}
\setkeys{Gin}{width=\linewidth,totalheight=\textheight,keepaspectratio}
\graphicspath{{graphics/}}

\title[The Descent of Artificial Intelligence]{The Descent of Artificial Intelligence: A Deep History of an Idea 400 Years in the Making} % longtitle shouldn't be necessary

% The following package makes prettier tables.  We're all about the bling!
\usepackage{booktabs}

% The units package provides nice, non-stacked fractions and better spacing
% for units.
\usepackage{units}

% The fancyvrb package lets us customize the formatting of verbatim
% environments.  We use a slightly smaller font.
\usepackage{fancyvrb}
\fvset{fontsize=\normalsize}

% Small sections of multiple columns
\usepackage{multicol}

% Provides paragraphs of dummy text
\usepackage{lipsum}

% These commands are used to pretty-print LaTeX commands
\newcommand{\doccmd}[1]{\texttt{\textbackslash#1}}% command name -- adds backslash automatically
\newcommand{\docopt}[1]{\ensuremath{\langle}\textrm{\textit{#1}}\ensuremath{\rangle}}% optional command argument
\newcommand{\docarg}[1]{\textrm{\textit{#1}}}% (required) command argument
\newenvironment{docspec}{\begin{quote}\noindent}{\end{quote}}% command specification environment
\newcommand{\docenv}[1]{\textsf{#1}}% environment name
\newcommand{\docpkg}[1]{\texttt{#1}}% package name
\newcommand{\doccls}[1]{\texttt{#1}}% document class name
\newcommand{\docclsopt}[1]{\texttt{#1}}% document class option name


\begin{document}

\begin{titlepage}

\begin{fullwidth}
\noindent\LARGE\emph{Book review
} \hspace{88mm}\includegraphics[height=1cm]{logo3.png}\\
\noindent\hrulefill\\
\vspace*{1em}
\noindent{\Huge{\emph{The Descent of Artificial Intelligence: A Deep History of an Idea 400 Years in the Making}\par}}

\vspace*{1.5em}

\noindent\LARGE{Sandra Braman}\par\marginnote{\emph{The Descent of Artificial Intelligence: A Deep History of an Idea 400 Years in the Making}, reviewed by Sandra Braman, \emph{History of Media Studies} 6 (2026), \href{https://doi.org/10.32376/d895a0ea.3bf4c334}{https://doi.org/10.32376/d895a0ea.3bf4c334}.} \vspace*{0.75em}
\vspace*{0.5em}
\noindent{{\large\emph{Michigan State University}, \href{mailto:bramansandra@gmail.com}{bramansandra@gmail.com}\par}} 

\end{fullwidth}

\vspace*{1em}

\noindent Kevin Padraic Donnelly\marginnote{\href{https://creativecommons.org/licenses/by-nc/4.0/}{\includegraphics[height=0.5cm]{by-nc.png}}}. \emph{The Descent of Artificial Intelligence: A Deep History of an Idea 400 Years in the Making}. 384 pp. University of Pittsburgh Press, 2024. \$45 (cloth).

\vspace*{2em}



\newthought{The word ``descent''} is in the title of Kevin Padraic Donnelly's book on
the history of artificial intelligence (AI) because, the author argues,
the ``achievement'' of AI intelligence has been the result of hundreds
of years of the development of social sciences that significantly
reduced our understanding of what it is to be human. Donnelly's deep
mining of the wide range of nuances in diverse views of what is meant by
``intelligence'' at all surrounds the reader with shifting prisms as to
where the boundaries of the human lie and possible relationships between
what we call the human and sources of knowledge and information that we
refer to as intelligence. The \emph{Oxford English Dictionary} confirms
what Donnelly is telling us. After the word ``intelligence'' first
appeared in print in the English language in 1390 in a poem by John
Gower, the history of its use, too, reflects a sense of it coming nearer
to being understood as the essence of being human, and then wandering
away from that sense. According to this book, the historical debate was
intense and raucous.

\emph{The Descent of Artificial Intelligence} is several books in one,
fractally replicating in structure the relations between the concept of
``intelligence'' and that of the human that are its subject matter. Some
of these are explicitly claimed by the book's author; others appear in
the course of reading. The technological edge of artificial intelligence
(AI)


\enlargethispage{2\baselineskip}

\vspace*{2em}

\noindent{\emph{History of Media Studies}, vol. 6, 2026}


\end{titlepage}

\noindent provides a wide heuristic umbrella, supporting those who begin
their discussions of AI by pointing out that it is not new.

Trying to achieve ``intelligence,'' however defined, has been a
Sisyphean effort; in author Kevin Padraic Donnelly's view, each approach
has not only failed but, in doing so, also reduced what the human is
understood to be. Even the social sciences, in this analysis, led to a
``new barbarism.'' Frustrated with the failed effort to achieve a common
understanding of intelligence itself, the turns to statistics and
probabilities were expected to solve the problem but did not. So much
for the social sciences. The author provides an extraordinary history of
thinking about human-machine relations, spelling out the wide range of
ways of conceptualizing the human, the machinic, and relations between
them as they have appeared across time and across societies.

From this perspective, the work is a sourcebook, each chapter
introducing a surprisingly large number of thinkers who have weighed in
on the issues. Thinkers with whose work many are generally familiar,
such as LaPlace and Rousseau, are discussed from perspectives that
broaden and deepen appreciation of their work as it was in conversation
with other authors in and before their times. Those of us who have
looked at the history of economic thought and of statistics will find
additional names familiar although, again, what is said about their
ideas took at least this reader in new directions. The significant
majority of writers analyzed by Donnelly, however, were new to me; I am
among those who will be treating the work as a sourcebook from now on.

Historians of the sociology of technology---and vendors---talk about
``unintended uses'' of technologies, but Donnelly reminds us there were
those who believed that it was only after a machine was made that people
began to understand its uses. As someone who has published on conceptual
distinctions among types of technologies and technological systems, this
book leaves me abashed, for the spectrum of types and dimensions along
which they can differ has multiplied and become far more complex many
times over. From this perspective, the book should be highly generative.
It would be quite valuable, for example, were an organizational
communication scholar to mine this book for insights, as the points of
intervention and structuration for group endeavors and inter-group
interactions will vary depending on how humans, machines, and their
relationships are understood. There are implications for many other
subfields of the media, communication, and information fields as well.

In the last decades of the twentieth century, it was fashionable not
only to rely upon ``postmodern'' thinkers such as Jean Baudrillard,
Michel Foucault, and Jacques Derrida, but to wonder why it was that so
many dominant figures in postmodern thought were French. Speculation at
the time was that it had to do with the fact that so many of them were
Catholic, with all of its technologies of belief. This work suggests
another causal factor; as \emph{The Descent of Artificial Intelligence}
tells us, French thinkers were probing relationships between symbols
(and the ``intelligence'' they convey) and the realities to which they
refer long before Baudrillard talked about hyperreality in language
associated with ``postmodernity.''

The scholarship is excellent; the writing doesn\textquotesingle t always
pass the same bar, overly trapped in the details. Moving densely through
several intricate but opposed or differentially nuanced arguments per
paragraph for much of the text, the reader wishes for an overall image
displaying the spectra of positions taken regarding the relationship of
intelligence to the human---whether it is immanent, whether the divine
is involved at all, whether it is essential to identity, whether
information is made or received or pushed. The writing is clear, but the
systematic presentations of argument, counter-argument, and twists on
arguments leave the reader more aware of trees than forest.

Despite this one weakness, the history of ideas offered here is
fascinating, deserving of study, paragraph by paragraph. It does not,
however, support the author's conclusion that there is nothing to fear
from artificial intelligence because the concept of intelligence has had
such a long and complex history. The realities of artificial
intelligence in today's environment, and future threats it may
pose---foreseeable and unforeseeable---remain real. Donnelly asserts
that technologies such as nuclear weapons and others are more dangerous
by far than AI. But the fear is that AI could drop nuclear bombs,
whether directly or indirectly, by misleading, confusing, or persuading
humans to do so. It is not an ``either/or'' regarding which technology
is most dangerous---it is ``and,'' with AI as multiplier of the
potential dangers of other technologies.

The comic character Philomena Cunk, created by Diana Morgan, reminds us
that thinking about thinking is the hardest thinking of all. As this
book demonstrates, it is.

\end{document}